\subsection{Eigenschaften der Verstärkerschaltungen}
\subsubsection{Aufgabenstellung}

Im Folgenden werden die Eigenschaften in hinsicht zur Eingagnsimpedanz und der
Phasenverschiebung bei invertierenden und nicht-invertierenden
Verstärkerschaltungen, wie in \autoref{fig:i_amplifier} und
\autoref{fig:ni_amplifier} dargestellt, verglichen.

\subsubsection{Eingangsimpedanz}

Der Eingangswiderstand ist gegeben durch das Verhältnis der Eingangsspannungsänderung zur Eingangsstromänderung:
\[
r_E = \frac{\partial u_E}{\partial i_E}
\]

\begin{wrapfigure}[7]{r}{0.4\textwidth}
	\centering\vspace{-11pt}
	\begin{circuitikz}[scale=0.8, transform shape]
		\myIA{A3}{(0,0)}{$R_1$}{$R_2$};
		\draw (A3-in) to[open, v=$U_E$, o-o] (A3-in |- A3-gnd) node[ground]{};
		\draw (A3-out) to[open, v=$U_A$, o-o] (A3-out |- A3-gnd) node[ground]{};
		\draw (A3-in) to[open, i=$I_E$] ++(0.2,0);
		\draw (A3-in) ++(0,-0.5) to[open, v=$U_1$] ++(2,0);
		\draw (A3-OA.-) to[open, v=$\SI{0}{\volt}$] (A3-OA.+);
	\end{circuitikz}
	\caption{Messschaltung i. Verstärker}\label{fig:i_amplifier_re}
\end{wrapfigure}

\paragraph{Invertierender Verstärker} Hier ist die Eingangsspanung $U_E$ aufgrund der virtuellen Masse am Knotenpunkt äquivalent zum Spannungsabfall $U_1$ über $R_1$. Der Eingagsstrom ist nun gegeben durch $I_E = \frac{U_1}{R_1}$, woraus folgt:
\[
r_E = \frac{U_E}{I_E} = \frac{U_1}{U_1/R_1} = R_1
\]

\begin{wrapfigure}[12]{r}{0.4\textwidth}
	\centering\vspace{-11pt}
	\begin{circuitikz}[scale=0.8, transform shape]
		\myNIA{A3}{(0,0)}{$R_1$}{$R_2$};
		\draw (A3-in) to[short, *-, i_<=$I_E$] ++(-1,0) coordinate(in) to[open, v=$U_E$, o-o] (in |- A3-gnd) node[ground]{};
		\draw (A3-in) to[R, l_=$r_D$] (A3-in |- A3-FB) -- (A3-FB);
		\draw (A3-out) to[open, v=$U_A$, o-o] (A3-out |- A3-gnd) node[ground]{};
	\end{circuitikz}
	\caption{Eingangswiderstand des nicht-invertierenden Verstärkers}
	\label{fig:ni_amplifier_re}
\end{wrapfigure}

\paragraph{Nicht-Invertierender Verstärker}

Der Eingangsstrom $I_E$ ist hier nur durch den Biasstrom $I_{B,P}$ des Operationsverstärkers bestimmt. Zieht man die Größen aus dem Modell des realen OPVs \cite[S.~66, Abb.~3.3]{EMTPraktikum2025} heran, ist zu sehen, dass der Differenzeingangswiderstand $r_D$ wie in \autoref{fig:ni_amplifier_re} in Serie geschalten ist, woraus folgt:
\[
r_E = R_1 + r_D \approx r_D \quad \text{da } r_D \gg R_1
\]
Auch die parallel leigenden Gleichtakteingangswiderstände $r_{GL}$ können
vernachlässigt werden, da dies ebenfalls sehr groß gegenüber $R_1$ sind.

\subsubsection{Phasenverschiebung}

\paragraph{Invertierender Verstärker} Das invertierende verhalten der Verstärkerschaltung führt, wie in \autoref{fig:bode_I} zu sehen, im Durchlassbereich des Frequenzgangs durch den vorzeichenwechsel der Spannung zu einer standardmäßigen Phasenverschiebung von $\varphi = \SI{180}{\degree}$

\paragraph{Nicht-Invertierender Verstärker} Hier bleibt das Ausgangssignal im durchlassbereich in Phase mit dem Eingangssignal.

\subsubsection{Erkenntnisse}

Es ist zu erkennen dass der nicht-invertierende Verstärker eine deutlich höhere
Eingangsimpedanz aufweist, da der Eingangsstrom nur durch den Biasstrom des
Operationsverstärkers bestimmt wird. Der invertierende Verstärker hat hingegen eine wesentlich geringer Eingangsimpedanz, die durch den Widerstand $R_1$ festgelegt ist.

Ein weiterer Unterschied liegt im Gleichtaktverhalten der beiden
Verstärkerschaltungen. Beim invertierenden Verstärker tritt keine Gleichtaktaussteuerung auf, da beide Eingänge des OPVs auf Massepotential liegen.
Beim nicht-invertierenden Verstärker hingegen wirkt die Spannung $U_1$ als Gleichtakteingangsspannung. Hat der OPV also eine beträchtliche Gleichtaktverstärkung (idealerweise $0$), muss das bei der Berechnung der Ausgangsspannung berücksichtigt werden.



